\subsection{更长时间策略的条件性交叉点}
\label{sec:schedule-amortization}

历史单桥诊断记录推进节省约 $\Delta M_0=0.119$ 秒，额外准备约 $\Delta P=0.206$--$0.240$ 秒。其计时窗口存在外部活动，被标记为 \texttt{diagnostic\_only}，且早于最终配置集合，不能据此产生实测最终加速比；但可以明确交叉点估计所需要的假设。

令 $N$ 为分段候选的总步数，包括未改变的前 $N_s=1024$ 步；其 $T=72$ 轨迹共有 $N_0=8192$ 步。若额外准备固定，推进节省仅随后段工作线性增长，则候选减去恒步控制的时间模型为
\begin{equation}
 \Delta T(N)=\Delta P-
       \frac{N-N_s}{N_0-N_s}\Delta M_0,\qquad
 N_* = N_s+(N_0-N_s)\frac{\Delta P}{\Delta M_0}.
 \label{eq:schedule-amortization}
\end{equation}
代入得候选步数 $N_*\simeq13{,}432$--$15{,}480$。若更粗略地按全部步数缩放节省，则得到 $N_*=N_0\Delta P/\Delta M_0\simeq14{,}181$--$16{,}522$。两者之差来自启动段的建模方式，不是新的测量。因此约 $1.5\times10^4$ 步只能作为有条件的实验规划估计，不能当作已观测阈值。模型没有证明 $T>72$ 时的精度、设备效率恒定，或相对另一合格方法的优势。多条短轨迹也只有在实际共享相关准备对象时才能摊销成本；仅缓存PTX并不共享空间矩阵或FFT计划。验证交叉点仍需要更长时域的资格检查和新隔离实验。
